\section{Notación principal}
\label{apx:notacion}

\begin{longtable}{P{3cm}P{11cm}}
\toprule
\textbf{Símbolo} & \textbf{Descripción} \\
\midrule
\endhead
$n$ & Número total de nodos del grafo, incluyendo el depósito. \\
$V$ & Conjunto de nodos del grafo, $V = \{1, 2, \dots, n\}$. \\
$V \setminus \{1\}$ & Conjunto de clientes (nodos distintos del depósito). \\
$E$ & Conjunto de aristas del grafo completo no dirigido. \\
$c_{ij}$ & Costo de recorrer la arista $\{i,j\}$. \\
$d_i$ & Demanda del cliente $i$. \\
$D$ & Capacidad homogénea de cada vehículo. \\
$m$ & Número de vehículos. \\
$\mathcal{R}$ & Conjunto de rutas de una solución. \\
$z$ & Costo total (distancia recorrida) de una solución. \\
\midrule
\multicolumn{2}{l}{\textit{Particle Swarm Optimization}} \\
\midrule
$N$ & Número de partículas del enjambre. \\
$d$ & Dimensión del espacio de búsqueda de PSO. \\
$\mathbf{x}_i$ & Posición de la partícula $i$ (solución candidata codificada). \\
$\mathbf{v}_i$ & Velocidad (desplazamiento) de la partícula $i$. \\
$\mathbf{p}_i$ & Mejor posición histórica de la partícula $i$ (\textit{pbest}). \\
$\mathbf{p}_g$ & Mejor posición del enjambre o del vecindario (\textit{gbest}). \\
$w$ & Peso de inercia. \\
$c_1,\, c_2$ & Coeficientes cognitivo y social. \\
$\mathbf{r}_1,\, \mathbf{r}_2$ & Vectores aleatorios $\sim U(0,1)^d$. \\
$\otimes$ & Producto de Hadamard (componente a componente). \\
$\chi$ & Factor de constricción. \\
\bottomrule
\end{longtable}

\section{Instancias de referencia}
\label{apx:instancias}

Las familias A, B y E de CVRPLIB se identifican mediante nombres del tipo
\texttt{A-n32-k5}, donde la letra indica la familia, \texttt{n} corresponde al número
total de nodos incluyendo depósito y \texttt{k} al número de vehículos de la solución de
referencia. Esta convención facilita la trazabilidad experimental y la comparación entre
estudios.

\section{Código implementado}
\label{apx:codigo}

% TODO: reemplazar por la URL del repositorio una vez publicado.
El código completo (módulo compartido \texttt{cvrp\_common.py} y los tres notebooks de
parametrización, experimentación y análisis) está disponible en:
\texttt{[TODO: enlace al repositorio]}.

A continuación se incluyen dos fragmentos representativos de la implementación de
referencia en Python puro (la versión efectivamente usada en las corridas está acelerada
con Numba, Sección~\ref{sec:parametrizacion}, pero es funcionalmente idéntica y fue
validada contra esta).

\begin{lstlisting}[language=Python, caption={Decodificación SR-1: lista de prioridad de clientes (SPV) e inserción de menor costo.}, basicstyle=\ttfamily\scriptsize, breaklines=true]
def decode_sr1(x, instance, n_vehicles, apply_local_search=True):
    n = instance.n_customers
    m = n_vehicles
    dist = instance.dist
    demands = instance.demands

    priority_list = _priority_list(x[:n])  # SPV: argsort ascendente + 1
    vehicle_priority = _vehicle_priority_matrix(x[n:n + 2*m], instance.coords, m)

    routes = [[] for _ in range(m)]
    loads = np.zeros(m)
    unassigned = []

    for customer in priority_list:
        demand = demands[customer]
        placed = False
        for j in vehicle_priority[customer - 1]:
            if loads[j] + demand > instance.capacity:
                continue
            pos, _delta = _best_insertion_cost(routes[j], customer, dist)
            routes[j].insert(pos, customer)
            loads[j] += demand
            if apply_local_search:
                routes[j] = _two_opt(routes[j], dist)
            placed = True
            break
        if not placed:
            unassigned.append(customer)  # penalizado en evaluate_particle
    return routes, unassigned
\end{lstlisting}

\begin{lstlisting}[language=Python, caption={Actualización de velocidad y posición del enjambre (Ecuaciones~\ref{eq:m2_velocidad}-\ref{eq:m2_posicion}).}, basicstyle=\ttfamily\scriptsize, breaklines=true]
w = params.w0 + (t / max(params.n_iter - 1, 1)) * (params.wf - params.w0)
r1 = rng.uniform(0, 1, size=(N, d))
r2 = rng.uniform(0, 1, size=(N, d))
V = w * V + params.c1 * r1 * (P - X) + params.c2 * r2 * (g[None, :] - X)
V = np.clip(V, -v_max, v_max)
X = X + V
below, above = X < x_min, X > x_max
X = np.clip(X, x_min, x_max)
V[below | above] = 0.0  # reflejo de velocidad a 0 en el borde
\end{lstlisting}
